% GENERATED FILE. Edit canonical lessons, then run npm run generate:books.
% canonical-source-hash: fnv1a64:a028a5b592c8680c
% canonical-lessons: GE-C407-rasen, GE-C407-blitzer, GE-C407-erwischen, GE-C407-bussgeld, GE-C407-strafzettel

\chapter{Zu schnell}
\label{ch:ge-zu-schnell}
\hlchaptermodality{407}

\begin{chapteropening}
\textbf{By the end of this chapter:} \emph{I can talk about speeding: racing along, the speed camera, getting caught, the fine and the penalty ticket.}

Complete the last lesson of chapter 407: I can talk about speeding: racing along, the speed camera, getting caught, the fine and the penalty ticket.
\end{chapteropening}

\section[rasen]{rasen — to speed, to race along}

\label{lesson:GE-C407-rasen}



\begin{quote}
Before the new one: say the German for the multi-storey car park, then the German for the parking ticket.
\end{quote}

\subsection*{You'll want to know: rasen}
\textbf{rasen} — \textquotedblleft{}to speed, to race along\textquotedblright{}.

\begin{quote}
Auf der Landstraße rast ein junger Mann an uns vorbei.
\end{quote}

\begin{quote}
\emph{On the country road a young man races past us.}
\end{quote}

Perfekt with \emph{sein} when someone moves fast: \emph{er ist gerast}.

\begin{cousinweb}
\emph{rasen} is an old Germanic verb that first meant 'to rage, to be wild' (EtymDB lists 'to rush, to rage'); its deeper origin is not settled.
\end{cousinweb}

\begin{grammarlens}[title={What you've built}]
One more word for telling, explaining and describing at B1.
\end{grammarlens}

\subsection*{Your turn}
\begin{itemize}
\raggedright
  \item \emph{Say it:} \emph{rasen}
  \item \emph{Say it:} \emph{rasen}, once more
  \item \emph{Say it:} a sentence of your own with \emph{rasen}
  \item \emph{Recall:} say the German for the multi-storey car park, then the German for the parking ticket, then say \emph{rasen} again
\end{itemize}

\begin{tcolorbox}[breakable,colback=teal!4,colframe=teal!35!black,title={Before you move on}]
What does rasen mean? (\textquotedblleft{}To speed, to race along\textquotedblright{}.) Say it once more.
\end{tcolorbox}
\section[der Blitzer]{der Blitzer — the speed camera (colloquial)}

\label{lesson:GE-C407-blitzer}



\begin{quote}
Before the new one: say the German for the parking ticket, then the German for to speed.
\end{quote}

\subsection*{You'll want to know: der Blitzer}
\textbf{der Blitzer} — \textquotedblleft{}the speed camera (colloquial)\textquotedblright{}.

\begin{quote}
Vor der Schule steht ein Blitzer.
\end{quote}

\begin{quote}
\emph{There is a speed camera in front of the school.}
\end{quote}

Same form in the plural. Everyday speech; officially \emph{die Geschwindigkeitskontrolle}.

\begin{cousinweb}
\emph{Blitzer} is built on \emph{blitzen} 'to flash', from \emph{Blitz} 'lightning', inherited from Middle High German \emph{blitze} (EtymDB): the camera flashes when a car is too fast.
\end{cousinweb}

\begin{grammarlens}[title={What you've built}]
One more word for telling, explaining and describing at B1.
\end{grammarlens}

\subsection*{Your turn}
\begin{itemize}
\raggedright
  \item \emph{Say it:} \emph{der Blitzer}
  \item \emph{Say it:} \emph{der Blitzer}, once more
  \item \emph{Say it:} a sentence of your own with \emph{der Blitzer}
  \item \emph{Recall:} say the German for the parking ticket, then the German for to speed, then say \emph{der Blitzer} again
\end{itemize}

\begin{tcolorbox}[breakable,colback=teal!4,colframe=teal!35!black,title={Before you move on}]
What does der Blitzer mean? (\textquotedblleft{}The speed camera (colloquial)\textquotedblright{}.) Say it once more.
\end{tcolorbox}
\section[erwischen]{erwischen — to catch (someone at something)}

\label{lesson:GE-C407-erwischen}



\begin{quote}
Before the new one: say the German for to speed, then the German for the speed camera.
\end{quote}

\subsection*{You'll want to know: erwischen}
\textbf{erwischen} — \textquotedblleft{}to catch (someone at something)\textquotedblright{}.

\begin{quote}
Der Blitzer hat ihn mit 80 erwischt.
\end{quote}

\begin{quote}
\emph{The speed camera caught him doing 80.}
\end{quote}

Everyday word. \emph{Den Bus gerade noch erwischen} — to manage to catch the bus.

\begin{cousinweb}
\emph{erwischen} is usually explained as built from \emph{er-} and \emph{wischen} 'to wipe, to make a quick movement', inherited from Middle High German \emph{wischen} (EtymDB lists \emph{wischen}): first a quick grab.
\end{cousinweb}

\begin{grammarlens}[title={What you've built}]
One more word for telling, explaining and describing at B1.
\end{grammarlens}

\subsection*{Your turn}
\begin{itemize}
\raggedright
  \item \emph{Say it:} \emph{erwischen}
  \item \emph{Say it:} \emph{erwischen}, once more
  \item \emph{Say it:} a sentence of your own with \emph{erwischen}
  \item \emph{Recall:} say the German for to speed, then the German for the speed camera, then say \emph{erwischen} again
\end{itemize}

\begin{tcolorbox}[breakable,colback=teal!4,colframe=teal!35!black,title={Before you move on}]
What does erwischen mean? (\textquotedblleft{}To catch (someone at something)\textquotedblright{}.) Say it once more.
\end{tcolorbox}
\section[das Bußgeld]{das Bußgeld — the fine (for a traffic offence)}

\label{lesson:GE-C407-bussgeld}



\begin{quote}
Before the new one: say the German for the speed camera, then the German for to catch.
\end{quote}

\subsection*{You'll want to know: das Bußgeld}
\textbf{das Bußgeld} — \textquotedblleft{}the fine (for a traffic offence)\textquotedblright{}.

\begin{quote}
Das Bußgeld für zu schnelles Fahren ist hoch.
\end{quote}

\begin{quote}
\emph{The fine for driving too fast is high.}
\end{quote}

Plural: \emph{die Bußgelder}. \emph{Der Bußgeldbescheid} is the letter that comes in the post.

\begin{cousinweb}
\emph{Bußgeld} is built from \emph{Buße} 'amends, penance', inherited from Old High German \emph{buoza} (EtymDB), and \emph{Geld} 'money', also an inherited Germanic word.
\end{cousinweb}

\begin{grammarlens}[title={What you've built}]
One more word for telling, explaining and describing at B1.
\end{grammarlens}

\subsection*{Your turn}
\begin{itemize}
\raggedright
  \item \emph{Say it:} \emph{das Bußgeld}
  \item \emph{Say it:} \emph{das Bußgeld}, once more
  \item \emph{Say it:} a sentence of your own with \emph{das Bußgeld}
  \item \emph{Recall:} say the German for the speed camera, then the German for to catch, then say \emph{das Bußgeld} again
\end{itemize}

\begin{tcolorbox}[breakable,colback=teal!4,colframe=teal!35!black,title={Before you move on}]
What does das Bußgeld mean? (\textquotedblleft{}The fine (for a traffic offence)\textquotedblright{}.) Say it once more.
\end{tcolorbox}
\section[der Strafzettel]{der Strafzettel — the penalty ticket, the fine slip}

\label{lesson:GE-C407-strafzettel}



\begin{quote}
Before the new one: say the German for to catch, then the German for the fine.
\end{quote}

\subsection*{You'll want to know: der Strafzettel}
\textbf{der Strafzettel} — \textquotedblleft{}the penalty ticket, the fine slip\textquotedblright{}.

\begin{quote}
Am Abend findet sie einen Strafzettel am Auto.
\end{quote}

\begin{quote}
\emph{In the evening she finds a ticket on the car.}
\end{quote}

Same form in the plural. In everyday speech also \emph{das Knöllchen}.

\begin{cousinweb}
\emph{Strafzettel} is built from \emph{Strafe} 'punishment, fine', inherited from Middle High German \emph{strāfe} (EtymDB), and \emph{Zettel} 'slip of paper', which goes back to Medieval Latin \emph{cedula} (EtymDB).
\end{cousinweb}

\begin{grammarlens}[title={What you've built}]
One more word for telling, explaining and describing at B1.
\end{grammarlens}

\subsection*{Your turn}
\begin{itemize}
\raggedright
  \item \emph{Say it:} \emph{der Strafzettel}
  \item \emph{Say it:} \emph{der Strafzettel}, once more
  \item \emph{Say it:} a sentence of your own with \emph{der Strafzettel}
  \item \emph{Recall:} say the German for to catch, then the German for the fine, then say \emph{der Strafzettel} again
\end{itemize}

\begin{tcolorbox}[breakable,colback=teal!4,colframe=teal!35!black,title={Before you move on}]
What does der Strafzettel mean? (\textquotedblleft{}The penalty ticket, the fine slip\textquotedblright{}.) Say it once more.
\end{tcolorbox}
